\begin{lemma}[The integral degree-two kernel of a cover]
\label{lem:B-Gm-proper-kernel}
\status{Proved without a Tate hypothesis}
Let \(q:T\to S\) be a morphism of finite-type
\(\F_p\)-schemes whose perfection is a \(v\)-cover,
and put \(T_a=T\times_S\cdots\times_ST\)
with \(a+1\) factors.  Then
\[
 \ker\bigl(H^2_{\et}(S,\Gm)[1/p]
       \longrightarrow H^2_{\et}(T,\Gm)[1/p]\bigr)
\]
is a finitely generated \(\Z[1/p]\)-module.  Consequently,
for every prime \(\ell\ne p\), the same kernel after tensoring
with \(\Z_{(\ell)}\) is finitely generated over
\(\Z_{(\ell)}\).  These are assertions about the full
cohomology groups, including their possible non-torsion parts.
In particular the lemma applies to proper surjections and to
finite-type \(h\)-covers obtained by combining a finite open
cover with proper surjections.
\end{lemma}
\begin{proof}
Write \(U^{\mathrm{perf}}\) for the perfection of a
characteristic-\(p\) scheme.  The universal homeomorphism
\(U^{\mathrm{perf}}\to U\) identifies the small \(\et\)
topoi.  Under this identification its multiplicative sheaf is
\(\varinjlim(\Gm\xrightarrow{p}\Gm\xrightarrow{p}\cdots)\):
this follows on affine \(\et\) objects by taking units of the
perfected ring.  Cohomology of a quasi-compact quasi-separated
scheme commutes with filtered colimits.  Therefore, naturally
in every finite-type \(U\),
\[
 H^i_{\et}(U^{\mathrm{perf}},\Gm)
       =H^i_{\et}(U,\Gm)[1/p].                     \tag{B.24f}
\]
The two inputs are respectively the Stacks Project's
Theorem~59.45.2 and Proposition~59.45.4, and Lemma~69.5.1
\cite{BStacksPerfectEtale,BStacksCohomColimits}.

On the category of perfect quasi-compact quasi-separated
schemes, use the \(v\)-topology.  Bhatt--Scholze's
Theorem~4.1(ii), with \(n=1\), says that vector bundles form
a \(v\)-stack; hence so does the groupoid of line bundles.
Its automorphism sheaf at the trivial bundle is \(\Gm\),
and it is the classifying stack of \(\Gm\) because line
bundles are already Zariski locally trivial.  Thus
\[
 H^0_v(U^{\mathrm{perf}},\Gm)
       =\Gamma(U^{\mathrm{perf}},\mathcal O^\times),
 \qquad
 H^1_v(U^{\mathrm{perf}},\Gm)
       =\operatorname{Pic}(U^{\mathrm{perf}}).
\]
For the change-of-topology morphism
\(\epsilon:(U^{\mathrm{perf}})_v\to
(U^{\mathrm{perf}})_{\et}\), this gives
\(\epsilon_*\Gm=\Gm\) and \(R^1\epsilon_*\Gm=0\),
since line bundles become trivial on Zariski covers.
The low-degree Leray sequence consequently supplies a natural
injection
\[
 H^2_{\et}(U^{\mathrm{perf}},\Gm)
       \lhook\joinrel\longrightarrow
 H^2_v(U^{\mathrm{perf}},\Gm).                     \tag{B.24g}
\]
No comparison in all higher degrees is being presumed.
Perfection commutes with the fibre products here, and the
perfected morphism is a \(v\)-cover by hypothesis.
A proper surjection becomes a \(v\)-cover, as do open
covers and their composites; the proper-surjection case and
its use in descent are also recorded in
Bhatt--Scholze's Proposition~4.5 and its proof
\cite[Theorem~4.1(ii), Proposition~4.5, author-version
pp.~15--17]{BBhattScholze2017}.

The ordinary cohomological descent spectral sequence for that
cover is
\[
 E_1^{a,b}=H^b_v(T_a^{\mathrm{perf}},\Gm)
       \Longrightarrow H^{a+b}_v(S^{\mathrm{perf}},\Gm).
\]
The kernel of restriction in total degree two has a two-step
filtration whose associated groups are \(E_\infty^{1,1}\)
and \(E_\infty^{2,0}\).  They are subquotients of
\[
 \operatorname{Pic}(T_1^{\mathrm{perf}})
       =\operatorname{Pic}(T_1)[1/p],\qquad
 \Gamma(T_2^{\mathrm{perf}},\mathcal O^\times)
       =\Gamma(T_2,\mathcal O^\times)[1/p].
\]
Both are finitely generated over \(\Z[1/p]\).
For the first this is Jaffe's Theorem~10.11.  For the second,
his Proposition~10.8 gives finite generation on the reduction,
and the kernel of reduction on units is killed by a bounded
power of \(p\).  Even if the map on global units is not
surjective, its image is a subgroup of a finitely generated
abelian group, which suffices
\cite[Proposition~10.8 and Theorem~10.11,
preprint pp.~42--43]{BJaffe1994}.
The ring \(\Z[1/p]\) is noetherian, so the indicated
subquotients, their extension, and all its submodules are
finitely generated.  By the natural injection (B.24g), the
kernel in the statement embeds in this \(v\)-cohomological
kernel.  This proves the assertion.  Further localization is
exact and proves the \(\ell\)-local assertion.
\end{proof}

\begin{proposition}[Proper singular bases introduce no further prime-to-\(p\) gate]
\label{prop:B-singular-proper-Brauer-gate}
\status{Proved equivalence; the surface Tate assertion remains unresolved}
Fix \(p\ne\ell\).  Divisor Tate for all smooth projective
surfaces over finite extensions of \(\F_p\) is equivalent to
the following assertion:
\[
 \begin{gathered}
 H^2_{\et}(S,\Gm)\otimes\Z_{(\ell)}
 \quad\hbox{is a finitely generated }\Z_{(\ell)}\hbox{-module}\\
 \hbox{for every proper finite-type }\F_p\hbox{-scheme }S.
 \end{gathered}
                                                               \tag{B.24h}
\]
Under these equivalent hypotheses, \(\Br(S)[\ell^\infty]\)
is finite for every such \(S\), including singular and
nonreduced ones.  In particular, if \(S\) is proper and
\(f:X\to S\) is a proper curve family, the \(\ell\)-primary
cokernel of \(\Br(S)\to\Br(X)\) is finite, whether or not
there is a global positive-degree line bundle on that family.
\end{proposition}
\begin{proof}
Assume surface Tate.  By Chow's lemma followed by de Jong's
projective alterations, choose a proper surjection \(T\to S\)
with \(T\) a finite disjoint union of smooth projective
varieties over finite fields.  Nilpotents can first be removed,
and the finitely many irreducible components treated separately
\cite[Theorem~4.1 and Remark~4.2]{c_dejong1996}.
The Kahn--Morrow reduction used in the proof of
Theorem~\ref{thm:B-relative-surface-gate} gives divisor Tate
on each component of \(T\).  Properness makes the arithmetic
\(H^1\) term in (B.24b) vanish by its weight \(-1\);
Kummer and finite cotype therefore make
\(\Br(T)[\ell^\infty]\) finite.  Since \(T\) is regular,
\(H^2_{\et}(T,\Gm)\) is torsion
\cite[Proposition~1.4 and Corollary~1.8]{BGrothendieckBrauerII1966}.
Consequently
\(H^2_{\et}(T,\Gm)\otimes\Z_{(\ell)}\) is finite.
Lemma~\ref{lem:B-Gm-proper-kernel} makes the kernel of
restriction from \(S\) finitely generated, so its source is
an extension of a finite module by a finitely generated module.
This proves (B.24h).

The torsion submodule of a finitely generated
\(\Z_{(\ell)}\)-module is a finite \(\ell\)-group.
Localization identifies the original \(\ell\)-primary
torsion with that submodule.  This proves the stated Brauer
finiteness; applying it directly to the proper total scheme
\(X\) proves the relative consequence.  It avoids replacing
the full Leray base term by its torsion subgroup.
Conversely, apply (B.24h) to a smooth projective surface.
Its full \(H^2(\Gm)\) is torsion, so the localized
module is finite.  The proper rational Kummer argument in
Theorem~\ref{thm:B-relative-surface-gate} then gives divisor
Tate.  This proves the equivalence.
\end{proof}

Properness is essential to this particular argument.  For a
smooth nonproper alteration \(T\), the group
\(\Br(T)[\ell^\infty]\) can contain
\(\Q_\ell/\Z_\ell\).  An extension of that group by a
finitely generated lattice need not be finitely generated:
\(\Z_{(\ell)}\subset\Q\) has quotient
\(\Q_\ell/\Z_\ell\) as an abstract abelian group.
Thus the kernel lemma alone does not control the nonproper
case.  The unit-cup argument in
Theorem~\ref{thm:B-all-bases-Brauer-units} supplies the
missing control and extends the relative conclusion to
every finite-type base.
